\section{THE DISCELL MODEL}
\label{sec:method}

We present the model in the order in which its parts depend on one another: the observed quantities and the spatial graph (\cref{sec:setting}); the generative model and the assumptions behind it (\cref{sec:generative}); the inference networks (\cref{sec:inference}); the objective (\cref{sec:objective}); the conditional-invariance regulariser and the held-out probe that grades it (\cref{sec:invariance}); the frozen-neighbour batching that makes the model trainable (\cref{sec:batching}); model selection (\cref{sec:selection}); and what is and is not identified, which is what turns a fit into a finding through the leakage sweep (\cref{sec:sweep}). Each design decision is stated with its reason; the longer arguments are in \cref{app:rationale} and the derivations in \cref{app:derivations}.

\subsection{Setting and Notation}
\label{sec:setting}

We consider a single tissue section profiled with a targeted in situ assay. For each of $N$ segmented cells $i$ we observe a count vector $\vx_i \in \Nat^G$ over $G$ genes, its total $\ell_i = \sum_g x_{ig}$, a cell-type label $t_i \in \{1,\dots,K\}$, and the morphology image around the cell. The label may be a curated annotation or a frozen unsupervised clustering and is treated as data, with two caveats that recur below: it was itself derived from the same contaminated counts, and its granularity is a modelling choice that sets what the decomposition is relative to (\cref{sec:invariance}). The total $\ell_i$ is conditioned on throughout and never modelled.

\paragraph{Spatial graph.}
Cells are nodes; two cells are adjacent if and only if their Voronoi cells share a face, i.e.\ no third cell lies between them (Delaunay adjacency). This choice has no neighbourhood count $k$ and no radius, adapts to local packing with an average degree near six, and uses only cell centroids, so it is insensitive to segmentation-boundary noise and to the nuclear-expansion segmentation that in situ platforms fall back on. Delaunay adjacency also connects cells across lumens, vessels and tears, so edges longer than a cut-off ($40\um$ here, well above the 95th percentile of edge length) are pruned. We write $\nbr{i}$ for the pruned neighbour set of $i$, which excludes $i$ itself; all its members are first-order neighbours, so no cell-size correction of distances is needed. The same graph serves every role below.

For each edge we record the centroid distance $d_{ij}$ and the length $f_{ij}$ of the shared Voronoi face, measured after clipping each Voronoi cell to a $30\um$ disc around its centroid so that faces at tissue borders and across gaps are bounded. The face length is a proxy for contact extent that is defined whether or not the segmented polygons touch. We deliberately do not use the length of shared segmented boundary: under nuclear-expansion segmentation, whether two polygons touch is largely a measurement of local density, and local density is a property of the niche that we want to keep out of the \emph{technical} part of the model, the part that describes measurement rather than biology (\cref{app:graph}).

\paragraph{Derived observables.}
From the graph and labels we form the neighbour-type composition $\vy_i \in \simplex{K-1}$, $y_{ik} = |\{j \in \nbr{i} : t_j = k\}| / |\nbr{i}|$, which excludes the cell itself. From the image we form an embedding $\vPhi_i \in \R^{\dPhi}$ of a patch centred on the cell, produced by a frozen pretrained vision model, in which a disc covering the cell has been masked out, so that $\vPhi_i$ describes the surrounding tissue and not the cell (\cref{app:image}). The two descriptors act at different scales: the graph is one hop, a few to thirty micrometres, while the masked patch describes tissue beyond roughly $25\um$; both are what we mean by ``the niche''. Here $\dPhi = 384$, the embedding's native width, since an ablation found no benefit from projecting it (\cref{app:calibration}). Three lookup tables are fixed before training: $\ybar(t)$, the mean composition of cells of type $t$; $\Phibar(t)$, the mean image-niche membership of type $t$; and $\bar{\vv}(t)$, the mean of the invariance block of \cref{sec:invariance} for type $t$. Isolated cells, those with no neighbours after pruning, have an empty composition and are excluded from all three tables.

\paragraph{Leakage kernel.}
The leakage kernel is a sparse matrix over the same edges, normalised over the neighbours of each connected cell,
\begin{equation}
  \beta_{ij} \;\propto\; f_{ij}\,\exp(-d_{ij}/\tau), \qquad \sum_{j \in \nbr{i}} \beta_{ij} = 1,
  \label{eq:beta}
\end{equation}
with contamination decay length $\tau = 20\um$; it is row-stochastic on connected cells, and an isolated cell keeps an all-zero row. It is precomputed once and never learned. At this cell spacing the decay is nearly inert and the kernel is dominated by the face lengths (\cref{app:pruning}); it is retained for platforms with sparser cells. Continuous geometry is allowed here, and only here: $\beta$ is not part of any conditioning variable and does not touch the estimand (\cref{app:graph}).

\subsection{Generative Model}
\label{sec:generative}

Each cell has two latent variables. The intrinsic state $\vz_i \in \R^{\dz}$ carries what the cell is; the response $\vw_i \in \R^{\dw}$ carries what its environment did to it. The environment itself enters as a deterministic context descriptor $\vc_i$ (defined in \cref{sec:inference}) that is shared by co-located cells. With $\mpsi$ a learned prior-mean network, $a$ a learned decoder, $\vB \in \R^{G \times \dw}$ a learned loading matrix and $\kappa \in [0,1)$ the fixed leak fraction,
\begin{align}
  \vz_i &\sim \N(\mathbf{0}, \vI), \label{eq:gen-z}\\
  \vw_i \mid \vc_i, t_i &\sim \N\big(\mpsi(\vc_i, t_i),\, \sigma_w^2 \vI\big), \label{eq:gen-w}\\
  \vrho_i &= \softmax\big(a(\vz_i) + \vB\vw_i\big), \label{eq:gen-rho}\\
  \rhobar_i &= \textstyle\sum_{j \in \nbr{i}} \beta_{ij}\, \vrho_j, \label{eq:gen-rhobar}\\
  \vp_i &= (1-\kappa)\, \vrho_i + \kappa\, \rhobar_i, \label{eq:gen-p}\\
  \vx_i \mid \ell_i &\sim \Mult(\ell_i, \vp_i). \label{eq:gen-x}
\end{align}
Here $\vrho_i \in \simplex{G-1}$ is the cell's own \emph{clean} composition, $\rhobar_i$ is the \emph{foreign influx}, the leakage-weighted average of its neighbours' clean compositions, and $\vp_i$ is the fitted multinomial probability. The decoder $a(\vz_i) \in \R^G$ is the cell's baseline programme in log space; column $k$ of $\vB$ is a gene programme and $w_{ik}$ is how far cell $i$ has moved along it. The prior variance is fixed at $\sigma_w = 1$, deliberately (\cref{sec:inference}).

For an isolated cell $\rhobar_i = \mathbf{0}$, and the mixture in \cref{eq:gen-p} would sum to $1-\kappa$, charging the cell $-\ell_i\log(1-\kappa)$ nats for leakage it cannot receive. The mixture is therefore renormalised per row,
\begin{equation}
  \vp_i = \frac{(1-\kappa)\,\vrho_i + \kappa\,\rhobar_i}{(1-\kappa) + \kappa\sum_{j \in \nbr{i}}\beta_{ij}},
  \label{eq:renorm}
\end{equation}
so that an isolated cell degenerates to $\vp_i = \vrho_i$, i.e.\ $\kappa_i = 0$: no neighbours, no foreign source. For connected cells the denominator is one. Probabilities are floored at $10^{-8}$ inside the logarithm. Isolated cells keep an all-zero composition $\vy_i$ and are not excluded from the regulariser or the probes of \cref{sec:invariance}; they are about two in a thousand cells on the primary section ($899$ of $407{,}120$ after pruning).

\paragraph{What the leakage term assumes.}
The leak fraction $\kappa$ is section-wide, not per cell, and there is no ambient term. A cell's contamination in reality varies with its size, its library size and the abundance of its neighbours; that per-cell variation is absorbed by $\vz$ and $\vw$ and is not addressed by the sweep. Diffuse background is assumed negligible relative to leakage from adjacent cells \needsource. The mixture in \cref{eq:gen-p} is exact under two assumptions: that the sources are independent Poisson variables, whose sum is again Poisson, and that the rate of leaked-in transcripts scales with the receiving cell's own rate rather than with the donors', which is where the per-cell variation just mentioned is hidden. The multinomial then follows from conditioning on the observed total (below).

\paragraph{What the decomposition means.}
The split between $\vz$ and $\vw$ is defined by two requirements. First, within a type, the intrinsic state should carry nothing that is predictable from the niche: $\I(\vz_i;\, [\vy_i, \vPhi_i] \mid t_i) = 0$. In \cref{eq:gen-z} $\vz_i$ is a priori independent of everything, and $\vy_i, \vPhi_i$ are not modelled, so this is not a property of the generative model; it is a constraint imposed on the variational family in \cref{sec:invariance}, in the sense of posterior regularisation \citep{ganchev2010}. Second, the environment reaches the counts only through the response and through leakage: $\vx_i \perp \vc_i \mid \vz_i, \vw_i, \{\vrho_j\}_{j \in \nbr{i}}$, which holds by construction because the decoder of \cref{eq:gen-rho} takes no context. Precisely, what the first requirement removes from $\vz$ is variation predictable from the neighbour composition and the masked image given the type. It does not remove intrinsic variation that happens to be spatially organised within a type, nor variation carried by the neighbours' own intrinsic codes; on a carcinoma section, where subclonal structure is the dominant within-type spatial variation, this is the main case rather than a corner case, and it is why the decomposition is relative to $t$ (\cref{sec:invariance}).

\paragraph{Why a multinomial.}
If $x_{ig} \sim \Pois(\lambda_{ig})$ independently, then $\vx_i \mid \sum_g x_{ig} = \ell_i$ is $\Mult(\ell_i, \boldsymbol{\lambda}_i / \sum_g \lambda_{ig})$. Since $\ell_i$ is observed, the multinomial is the coherent likelihood: no library-size latent and no double counting of depth. Poisson sources are closed under addition, which is what makes the leakage mixture exact; negative-binomial sources with gene-specific means are not in general, and a negative binomial on the mixed counts would let the dispersion absorb part of what the mixture should explain. At the depth of in situ data, tens to a few hundred counts over thousands of genes, gene-wise overdispersion cannot be estimated in any case, and a low-dimensional amortised $\vz$ absorbs much of what a dispersion parameter captures in coarser models. Should posterior predictive checks show overdispersion at higher depth, the appropriate extension is a Dirichlet-multinomial with one concentration scalar, not $G$ dispersions.

\paragraph{Why $\kappa$ is fixed.}
Leakage and response differ in form: one adds a neighbour's own profile through a fixed contact kernel, the other shifts the cell's own programmes along learned directions. From the cell-by-gene count matrix alone, however, the split between them is pinned only by those parametric forms, and a decoder flexible enough to fit the data can mimic either, since both make a cell resemble its neighbours and neighbour composition predicts both. We do not trust the likelihood to choose $\kappa$. Instead we fix it, fit the model independently at every point of a grid, and report each readout as a function of $\kappa$ (\cref{sec:sweep}). $\kappa = 0$ is the nested no-correction null. The grid is anchored at zero and extends to $0.4$, close to the upper range of published misassignment estimates for this platform \needsource; a data-driven upper bound from an atlas-defined marker set, and the transcript-level information that would identify $\kappa$ from a nuclear/extranuclear split of the counts, are discussed in \cref{app:kappa}.

\paragraph{Why the decoder has no type input.}
The decoder in \cref{eq:gen-rho} takes $\vz_i$ and $\vw_i$ only, so the only route from identity to expression runs through $\vz$. If $\vz$ collapsed to a function of the label, both reconstruction terms in \cref{sec:objective} would lose all within-type variance, a large fitting penalty, and that is what holds $\vz$ open. With the label in the decoder, $\vz$ would become a residual-only code and would no longer be interpretable as identity. What $\vz$ carries beyond the label is what the label discretises away: cell-cycle phase, clonal identity, position on a differentiation continuum, stress and metabolic state, states acquired in a niche the cell has since left, and systemic signals that affect every cell of a type regardless of position. The last is environmental in the colloquial sense but not niche-predictable, and it belongs in $\vz$. Two remarks follow. Since the encoder of \cref{sec:inference} receives the label, the agreement between $\vz$ and $t$ measures whether $\vz$ \emph{retains} type, not whether it discovers it; its collapse is the detector for the marginal-invariance failure mode (\cref{sec:selection}). And coarse labels leave more for $\vz$ to carry; whether $\vz$ separates cell-cycle phases within a type is checked in \cref{sec:allegiance}.

\subsection{Inference Networks}
\label{sec:inference}

\Cref{tab:learned} summarises what is learned and what is fixed. All networks are multilayer perceptrons unless stated; widths are in \cref{tab:architecture}.

\begin{table}[t]
\centering
\caption{Learned and fixed quantities. $G = 5{,}101$, $K = 18$, $\dz = 20$, $\dw = 6$, $\dPhi = 384$ on the primary section.}
\label{tab:learned}
\small
\begin{tabular}{@{}p{0.47\columnwidth}p{0.47\columnwidth}@{}}
\toprule
Learned & Fixed before training \\
\midrule
encoder of $q(\vz_i \mid \vx_i, t_i)$; attention layer and type query $\embed(t)$; prior mean $\mpsi$; encoder of $q(\vw_i \mid \cdot)$; decoder $a$; loadings $\vB$; adversary heads (separate optimiser) &
graph, kernel $\beta$, $\tau$; leak fraction $\kappa$ (swept); prior scale $\sigma_w = 1$; image embedding $\vPhi$ and its principal components; image-niche memberships $\ephi$; tables $\ybar(t)$, $\Phibar(t)$, $\bar{\vv}(t)$; labels $t$; totals $\ell$ \\
\bottomrule
\end{tabular}
\end{table}

\paragraph{The context descriptor $\vc_i$.}
One vector per cell summarising who and what is nearby. It is deterministic, is shared by co-located cells, and is the conditioning variable of the model. It is computed by a single graph-attention layer \citep{brody2022} over the pruned graph, concatenated with the image embedding and an isolation flag:
\begin{align}
  \mathbf{h}_j &= \big[\onehot(t_j),\; \sg{\vmu_z(\vx_j, t_j)}\big], \label{eq:context-src}\\
  \mathbf{g}_i &= \GAT\big(\text{query} = \embed(t_i),\; \{\mathbf{h}_j\}_{j \in \nbr{i}}\big), \label{eq:context-gat}\\
  \vc_i &= \mathbf{g}_i \;\oplus\; \vPhi_i \;\oplus\; \mathbf{1}[\nbr{i} = \emptyset]. \label{eq:context}
\end{align}
Here $\vmu_z(\vx_j, t_j)$ is the posterior mean of $\vz_j$ (below), $\sg{\cdot}$ denotes a stop-gradient, and $\embed(t_i)$ is a learned per-type vector of the same width as the source features. The layer has no self-loops, no edge features, and its heads are averaged. Its attention weights $\att_{ij}$ satisfy $\sum_j \att_{ij} = 1$. Three decisions here are load-bearing.

\emph{The query is the type, not $\vz_i$.} In attention of this form the destination feature enters the attention logits but never the aggregated value, so using $\vz_i$ as the query would inject no ego content directly. It would, however, open a selection channel: attention can learn $\att_{ij} \propto \exp(\simop(\vz_i, \vz_j))$ and reconstruct the cell through its most similar neighbours, so that $\vc_i$ becomes a mirror of $\vz_i$. The incentive exists because an ego-informed prior on $\vw_i$ is free decoder capacity, and it is strongest in homotypic regions, where $\kappa$ is also least identified. The consequence would be that $\mpsi$ already knows the cell, the divergence $\KL(q(\vw_i)\,\|\,p(\vw_i \mid \vc_i, t_i))$ collapses, and $\vw_i$ measures deviation from \emph{self} rather than from the expected response. Querying with the type closes the selection channel at no cost to the prior, since $\mpsi(\vc_i, t_i)$ already receives $t_i$. A value channel remains: in homotypic regions the neighbours' codes resemble the cell's through shared lineage and through symmetric contamination, so $\vc_i$ can correlate with $\vz_i$ regardless of the query. That residual is monitored, not assumed away (\cref{sec:selection}; \cref{app:mirror}).

\emph{No per-edge geometry enters the context.} The graph says who is adjacent; the image describes how packed and how structured the region is; per-edge distances and contact extents are withheld from $\vc_i$. Three reasons, in decreasing order of weight. The concrete risk is collinearity: an attention layer given edge features can rediscover the leakage kernel $\beta$, after which $\vw$ becomes collinear with $\rhobar$ and the two mechanisms the sweep is meant to separate merge inside the model (\cref{app:graph}). Within a first-order neighbourhood the distances span perhaps $5$ to $30\um$, most of which is local density, which the image already carries. And the intended reading of the context is that neighbour-type composition is the exposure of interest while the image and the graph are conditioning covariates; a hypothetical change in composition is a meaningful question about a tissue in a way that ``the same neighbours, further away'' is not, and putting per-edge geometry into the conditioning variable would mix the two. The estimand is used only informally here; counterfactual machinery is future work (\cref{sec:conclusion}).

\emph{Isolated cells.} A cell with no in-edges has an empty attention set; the graph part of $\vc_i$ is set to zero and the flag is raised. The image part stays, since $\vPhi_i$ is measured data defined for every cell. Together with \cref{eq:renorm}, an isolated cell is modelled with $\kappa_i = 0$ and an image-only context.

\paragraph{Posterior over $\vz$.}
The posterior $q(\vz_i \mid \vx_i, t_i)$ is a diagonal Gaussian produced by a network whose input is the library-normalised, log-transformed count vector together with the log total and the one-hot type,
\begin{equation}
  \big[\log(1 + \vx_i\, \tilde\ell / \ell_i),\; \log \ell_i,\; \onehot(t_i)\big], \qquad \tilde\ell = \operatorname{median}_i \ell_i.
  \label{eq:encz-input}
\end{equation}
The encoder \emph{never sees} $\vc_i$: invariance by architecture, which the regulariser of \cref{sec:invariance} then completes. Library normalisation before the log matters because a raw $\log(1+\vx)$ entangles depth with composition, so that a 300-count and a 50-count cell of the same type look very different; since $\ell_i$ is conditioned on rather than modelled, $\vz$ should not spend capacity on it. Total counts are nonetheless weakly informative about state, so $\log \ell_i$ enters as one extra scalar rather than being discarded.

\paragraph{Prior and posterior over $\vw$.}
The prior is what a typical cell of this type would do in this context; the posterior is what this cell did:
\begin{align}
  p(\vw_i \mid \vc_i, t_i) &= \N\big(\mpsi(\vc_i, t_i),\, \sigma_w^2 \vI\big), \label{eq:prior-w}\\
  q(\vw_i \mid \vc_i, t_i, \vz_i, \vx_i) &= \N\big(\vmu_{w,i},\, \diag \boldsymbol{\nu}_{w,i}^2\big), \label{eq:post-w}
\end{align}
with $\boldsymbol{\nu}_{w,i}$ the learned posterior scale, as distinct from the fixed prior scale $\sigma_w$. Both distributions are functions of the context, but only the posterior sees the cell itself, through the sampled $\vz_i$ and through the encoded counts of \cref{eq:encz-input}; we abbreviate the posterior as $q(\vw_i)$ where its arguments are clear. The divergence $\KL(q \,\|\, p)$ reads as how far this cell deviates from its expected response, an anomaly score that the objective provides at no extra cost; whether it carries usable per-cell information at the operating point is an empirical question, reported in \cref{sec:w-degeneracy}. The prior is also what would make counterfactuals possible, by drawing $\vw' \sim p(\vw \mid \vc', t)$ and decoding, with abduction $\vw'_i = \mpsi(\vc', t_i) + \sigma_w \hat{\boldsymbol\varepsilon}_i$ carrying the cell's own residual $\hat{\boldsymbol\varepsilon}_i = (\vw_i - \mpsi(\vc_i, t_i))/\sigma_w$; that machinery is not exercised in this paper.

The direct path from $\vx_i$ into the posterior is necessary. The conditional independence $\vw \perp \vx \mid \vz, \vc$ is false by explaining-away: $\vx$ is generated from both $\vz$ and $\vw$, so knowing $\vx$ and $\vz$ informs $\vw$ directly, which is also why \citet{slavutsky2025} condition their condition-specific posterior on the data rather than on the shared code. In our setting $\vz$ is lossy in exactly the wrong direction: it is low-dimensional, the invariance regulariser removes niche-predictable variation from it, and that variation \emph{is} the spatial response that $\vw$ exists to capture. Without the direct path, $\vw_i$ collapses to $\E[\text{response} \mid \text{type}, \text{niche}]$, identical for co-located similar cells, and the divergence vanishes (\cref{app:qw}).

\paragraph{Learned and fixed variances.}
Both posteriors carry learned diagonal variances; this is what makes the reparameterised sample meaningful \citep{kingma2014,rezende2014} and lets uncertainty vary per cell, so that a 30-count cell can have a wider $q(\vz)$ than a 300-count cell. The prior scale $\sigma_w$ is fixed at one, for two reasons of which the first is decisive. A learned \emph{constant} $\sigma_w$ is a gauge freedom: $\vw$ and $\vB$ enter the decoder only as $\vB\vw_i$, so rescaling dimension $k$ of $\vw$ by $s$ and column $k$ of $\vB$ by $1/s$ leaves the likelihood unchanged, and a learned constant prior scale would absorb the compensating change in the prior; fixing it pins the scale. A learned \emph{conditional} $\sigma_w(\vc, t)$ would be meaningful, claiming that some niches produce more variable responses than others, but it opens a degenerate direction: the model can shrink $\KL(q\|p)$ for free by inflating $\sigma_w(\vc,t)$ wherever the penalty bites, making the prior vague rather than making $\vw$ predictable. If it is wanted, it must be bounded (\cref{app:variances}).

\paragraph{Decoder.}
The decoder is \cref{eq:gen-rho}: a network $a(\vz_i) \in \R^G$ plus the linear shift $\vB \vw_i$, followed by a softmax over genes. It takes no type input, for the reason given in \cref{sec:generative}. If reconstruction is poor the remedy is to widen $\vz$ or the network, not to add $t$.

\subsection{Objective}
\label{sec:objective}

With the neighbour quantities $\vc_i$ and $\rhobar_i$ held constant (\cref{sec:batching}), the evidence lower bound on $\log p(\vx_i \mid \vc_i, t_i)$ under the variational family $q(\vz_i)\, q(\vw_i \mid \vc_i, t_i, \vz_i, \vx_i)$ is
\begin{equation}
\begin{split}
  L_i = \E_{q(\vz_i)}\Big[&\E_{q(\vw_i \mid \vz_i)}\big[\log p(\vx_i \mid \vz_i, \vw_i)\big] \\
  &- \KL\big(q(\vw_i \mid \vz_i)\,\|\,p(\vw_i \mid \vc_i, t_i)\big)\Big] \\
  &- \KL\big(q(\vz_i)\,\|\,p(\vz_i)\big),
\end{split}
\label{eq:bound-a}
\end{equation}
where the $\vw$-divergence sits inside the outer expectation because $q(\vw_i \mid \cdot)$ conditions on the sampled $\vz_i$; evaluating the closed-form Gaussian divergence at the one reparameterised draw is an unbiased one-sample estimate of it (\cref{app:bound}). A second bound on the same data is obtained for the model in which $\vw$ is not inferred, i.e.\ $q(\vw_i) := p(\vw_i \mid \vc_i, t_i)$, which zeroes that divergence:
\begin{equation}
\begin{split}
  L_i^{z} = \E_{q(\vz_i)}\E_{p(\vw_i \mid \vc_i, t_i)}\big[\log p(\vx_i \mid \vz_i, \vw_i)\big] \\
  - \KL\big(q(\vz_i)\,\|\,p(\vz_i)\big).
\end{split}
\label{eq:bound-b}
\end{equation}
We call this the \emph{intrinsic path}, and it is load-bearing. The first reconstruction term does reward $\vz$ for carrying signal, since $a(\vz)$ is the only route to all $G$ genes, but with the counts also wired into $q(\vw)$ the allocation between $\vz$ and $\vw$ is ambiguous, and the second bound breaks the tie toward $\vz$: under it $\vz$ must explain the cell given only the \emph{expected} response. Because it runs through the same leakage mixture as the first bound, it never asserts $\kappa = 0$ and so never contradicts it (\cref{app:pathb}).

The training objective sums the two bounds, with a weight $\omega$ on the intrinsic path and per-term weights on the divergences,
\begin{equation}
\begin{split}
  J = \frac{1}{|\mathcal{S}|}\sum_{i \in \mathcal{S}} \Big[
  & \tfrac{1}{\ell_i}\textstyle\sum_g x_{ig} \log p_{ig}(\vz_i, \vw_i) \\
  & + \omega\, \tfrac{1}{\ell_i}\textstyle\sum_g x_{ig} \log p_{ig}(\vz_i, \wbreve_i) \\
  & - (1+\omega)\,\alphaz\, \KL\big(q(\vz_i)\,\|\,\N(\mathbf{0},\vI)\big) \\
  & - \alphaw\, \KL\big(q(\vw_i)\,\|\,\N(\mpsi(\vc_i,t_i), \sigma_w^2\vI)\big) \Big] \\
  & - \alphaa\, \Pen,
\end{split}
\label{eq:objective}
\end{equation}
where $\mathcal{S}$ is the set of cells whose loss is evaluated in a batch, $\wbreve_i \sim p(\vw_i \mid \vc_i, t_i)$ is a draw from the prior, and $\Pen$ is the invariance term of \cref{sec:invariance}, a functional of the whole batch; under the adversarial escalation it is replaced by the per-cell mean $\frac{1}{|\mathcal{S}|}\sum_i \Adv_i$. One decoder is evaluated twice, with $\vw$ from the posterior in the first reconstruction term and from the prior in the second. The factor $(1+\omega)$ is what keeps $J$ proportional to $L_i + \omega L_i^{z}$: each of \cref{eq:bound-a,eq:bound-b} carries its own $-\KL(q(\vz_i)\|p(\vz_i))$, so their weighted sum contains $(1+\omega)$ copies of it, and omitting the second copy, the natural mistake, breaks the correspondence even at unit weights.

\paragraph{Scaling.}
The reconstruction terms are divided by $\ell_i$. Unscaled, they are of order $-200$ to $-400$ nats per cell against divergences of order $1$ to $10$, and their magnitude varies severalfold between sections, so weights tuned on one section would not transfer to the next. The scaling changes what the weights mean: reconstruction becomes $O(1)$ while the divergences stay absolute, so $\alphaz = 1$ prices a latent's information roughly $\ell$-fold above the unscaled bound and collapses $\vz$. At $\alphaz = \alphaw = 1/\lbar$, with $\lbar$ the mean total count, $J$ is proportional to $L_i + \omega L_i^{z}$ for a cell of average depth, and that point, not $1$, is the centre of any sweep over the weights. A per-cell scaling against global weights also means that a low-count cell's divergences are weighted more heavily relative to its likelihood than a high-count cell's, pulling sparse cells toward the prior; we accept this as the price of transferable weights. Once any weight departs from $1/\lbar$, or through the per-cell scaling itself, $J$ is a weighted surrogate rather than a bound; and because $\vx_i$ appears in both reconstruction terms, $q$ is not the posterior of any single model even at the bound-equivalent weights. The objective is coherent as an M-estimation target, and uncertainty in the reported effects comes from the leakage sweep, not from $q$. The hyperparameters are therefore $\omega$, $\alphaz$, $\alphaw$ and the response width $\dw$, plus $\alphaa$ and, under escalation, the adversary's update schedule; $\kappa$ is not a hyperparameter but the sweep axis.

\subsection{Conditional Invariance of \texorpdfstring{$\vz$}{z}}
\label{sec:invariance}

Architecture alone does not keep the environment out of $\vz$. The counts $\vx_i$ carry environmental signal, and with $\dz > \dw$ the intrinsic latent will absorb it given the chance. The target is
\begin{equation}
  \I\big(\vz_i;\, [\vy_i, \vPhi_i] \mid t_i\big) = 0.
  \label{eq:invariance-target}
\end{equation}
Two features of this target are deliberate. It is \emph{conditional on type}. Types cluster in space, so $t_i$ and $\vy_i$ are strongly dependent, and a marginal constraint $\I(\vz; \vy) = 0$ deletes cell type from $\vz$, because carrying ``I am a T cell'' predicts ``my neighbours are T cells''. This is the documented failure of conditional subspace VAEs \citep{klys2018}, whose invariance drives the nuisance information to zero while collapsing the code's agreement with cell type \citep{slavutsky2025} \todo{verify the numbers quoted in the design notes against the DisCoVR paper}. The cost of conditioning is that intrinsic variation that is spatially organised within a type is pushed into $\vw$, tumour subclones being the sharp case; the decomposition is relative to $t$, not absolute, and coarse labels are preferable to fine ones. It also targets \emph{both} $\vy$ and $\vPhi$. If only composition were penalised, $\vz$ would be free to absorb the morphology-defined niche (vessel proximity, fibrosis, density, necrosis), and $\vw$ would under-report morphology-driven effects; $\vy$ and $\vPhi$ are correlated, so the part that matters is the residual that is image-predictable but not composition-predictable, and that residual is what would go unpenalised. Neither the penalty below nor the adversary is part of any bound; both are posterior regularisation \citep{ganchev2010}, a constraint on the variational family written as a penalty, which is the status the adversarial term also has in \citet{slavutsky2025} (\cref{app:posterior-reg}).

\paragraph{Closed-form penalty.}
The first regulariser is the Gaussian estimate of the conditional mutual information, which needs no second optimiser. With $\vv_i = [\vy_i^{(-K)}, \mathrm{PC}(\vPhi_i)]$, where $\vy^{(-K)}$ drops one composition column (the simplex constraint makes the full block singular) and $\mathrm{PC}(\vPhi)$ keeps the leading principal components of the image embedding,
\begin{equation}
\begin{split}
  \Pen = \sum_t p(t)\cdot \tfrac{1}{2}\Big[&\logdet \Sighat_{\vv \mid t} + \logdet \Sighat_{\vz \mid t} \\
  &- \logdet \Sighat_{[\vz,\vv] \mid t}\Big],
\end{split}
\label{eq:penalty}
\end{equation}
evaluated on the posterior means $\vmu_z$ rather than on samples, since encoder noise would dilute the measured dependence and hiding dependence behind noise is exactly what the penalty must not permit (\cref{app:gaussian-mi}). Three estimation details are not optional at realistic batch sizes. The image block is kept small (a dozen principal components), which keeps the joint block well conditioned and retains the components that carry the niche signal. The per-type moments are accumulated as exponential moving averages across batches (rate $0.05$, an effective window of about twenty batches), because a spatial batch contains few cells of any rare type; a type enters the penalty only once its window-weighted cell count reaches ten times the joint dimension, and the fraction of cells excluded by this support floor is reported, since those cells go unregularised. The type weights $p(t)$ are the section-wide type frequencies and are not renormalised after an exclusion. Regularisation shrinks toward the diagonal \citep{ledoit2004,schafer2005}, $\Sighat \leftarrow (1-\gamma)\Sighat + \gamma\,\diag(\Sighat)$, which damps only the correlations; an additive ridge $+\lambda\vI$ would inflate every variance by an amount that depends on how ill-conditioned the estimate was, biasing the log-determinant upward most for rare types and weakening the penalty exactly where data are thinnest. Two further details concern the gradient. The value of \cref{eq:penalty} is taken from the moving-average moments but its gradient from the current batch at full scale, a straight-through estimator \citep{bengio2013}; otherwise the only differentiable path runs through the averaging weight, and the effective strength of $\alphaa$ would silently scale with the averaging constant. And the log-determinants are computed on correlation rather than covariance matrices: mutual information is invariant to per-dimension rescaling, so the value is identical, while the gradient of a log-determinant, which involves the inverse matrix, stays bounded when a dimension has near-zero variance, as the composition of cells deep inside one niche does. A jitter of $10^{-5}\vI$ is added to the unit-diagonal correlation matrix before the factorisation; on a correlation matrix this is a uniform relative perturbation, not the variance-dependent ridge rejected above.

\paragraph{Adversarial escalation.}
The Gaussian penalty removes linear dependence only. When the probe below finds that dependence survives it, the penalty is replaced by an adversary \citep{louizos2016,xie2017,ganin2016}: two heads with parameters $\xi$, $\hat{\vy}_\xi, \hat{e}^{\Phi}_\xi : (\vmu_{z,i}, t_i) \mapsto \simplex{K-1}, \simplex{\Ephi-1}$, trained by a separate optimiser to predict the composition and an image-niche membership from the posterior mean of the intrinsic code (evaluated on the mean for the same reason as \cref{eq:penalty}), and an encoder term
\begin{equation}
\begin{split}
  \Adv_i = \;&\CE\big(\vy_i, \ybar(t_i)\big) - \CE\big(\vy_i, \hat{\vy}_\xi(\vmu_{z,i}, t_i)\big) \\
  +\;&\CE\big(\ephi_i, \Phibar(t_i)\big) - \CE\big(\ephi_i, \hat{e}^{\Phi}_\xi(\vmu_{z,i}, t_i)\big),
\end{split}
\label{eq:adv}
\end{equation}
which is the heads' excess predictive skill over knowing the type alone, in categorical cross-entropy. The adversary maximises it, the encoder minimises it, and it is zero at the optimum; the two halves are reported separately. The image target $\ephi_i \in \simplex{\Ephi-1}$ is a soft cluster membership of $\vPhi_i$ over $\Ephi$ image niches, obtained by $k$-means on its leading principal components (never on the raw high-dimensional embedding, where distances are noise-dominated), fitted once and frozen. We set $\Ephi = K$ so that the two heads face comparable difficulty and one weight $\alphaa$ serves both; soft rather than hard assignments give a smoother target with no boundary artefacts. The adversary and the probe below therefore target the image through different summaries, discrete memberships and continuous components respectively; the memberships are the natural target for a classification head, the components for the Gaussian penalty and its probe. A weak adversary is worse than none: heads that are updated too slowly are defeated by the encoder while an independent probe still finds the leak, so head strength must be verified by the probe, never by the heads' own training loss (\cref{app:calibration}). This is the observation of \citet{elazar2018} that adversarial removal is not enough without an independent check.

\paragraph{The probe, and how to decide.}
Whether the regulariser worked is decided by a \emph{fresh} predictor of the invariance block $\vv_i$ from $(\vmu_{z,i}, t_i)$, trained on a spatially separate training split and evaluated on held-out cells: never the training adversary, and never the training cells. We report the excess held-out skill over the type-only baseline,
\begin{equation}
  \dCE = \tfrac{1}{d_v}\Big(\E\big\|\vv - \bar{\vv}(t)\big\|^2 - \E\big\|\vv - \hat{\vv}_{\mathrm{probe}}(\vmu_z, t)\big\|^2\Big),
  \label{eq:probe}
\end{equation}
the Gaussian form of the cross-entropy gain on the continuous block (mean squared error over its $d_v = K - 1 + 12$ components, which is the cross-entropy up to constants), with two reference scales that make it interpretable: the \emph{noise floor}, the value obtained when $\vmu_z$ is permuted across cells within type, and the \emph{uncontrolled baseline}, the value from a model fit at $\alphaa = 0$. The probe is fitted in two families, a ridge regression and a small multilayer perceptron (\cref{tab:architecture}); the permutation floor is reported for the ridge probe, and the decision is taken on the nonlinear one. In our experiments the linear probe sits at the floor while the nonlinear probe still reads a substantial fraction of the uncontrolled leak (\cref{app:calibration}). The rule, as applied, is to escalate from \cref{eq:penalty} to \cref{eq:adv} when the nonlinear probe of a converged closed-form fit is clearly above the floor, and then to take the smallest $\alphaa$ that brings the nonlinear probe within a small fraction of the uncontrolled baseline without a loss in the agreement between $\vz$ and cell type or in held-out reconstruction. The agreement is the normalised mutual information \citep{strehl2002} between $t$ and a $k$-means clustering of $\vmu_z$ with $k = K$ clusters. If no such $\alphaa$ exists, the invariance is costing cell type and the label set is too fine. \todo{state the fractions actually used once the calibration scans are repeated}

\subsection{Frozen Neighbours and Batching}
\label{sec:batching}

Neighbour quantities are treated as data, not as things to back-propagate through. The neighbour codes $\vmu_z(\vx_j, t_j)$ in \cref{eq:context-src} and the neighbour compositions $\vrho_j$ in \cref{eq:gen-rhobar} both enter under a stop-gradient. The mechanism this prevents is a degeneracy rather than an explosion: with the path open, cell $j$'s clean composition is trained to absorb cell $i$'s residual through $\rhobar_i$, so $\vrho_j$ stops being $j$'s own profile, and the coupling feeds back through $\rhobar_j$. \citet{ergen2025} disable gradients on precisely this path after observing mode collapse with it enabled; a first-order feedback argument bounds the amplification through the coupling by $1/(1-\kappa)$ (\cref{app:amplification}). The neighbours' responses $\vw_j$ needed for $\vrho_j$ are posterior draws computed in the same batched pass and likewise placed under a stop-gradient.

A tempting shortcut is to compute the foreign influx from the baseline programme alone, $\vrho_j \approx \softmax(a(\vz_j))$, which avoids the neighbours' context entirely. It errs in the direction that matters most. Neighbours share niches, so $\vw_j \approx \vw_i$; dropping $\vw_j$ strips the local spatial programme out of $\rhobar_i$, and whatever the leakage term then fails to absorb is explained by $\vw_i$ instead, \emph{inflating} apparent spatial effects, the exact failure the leakage sweep exists to guard against.

\paragraph{Two hops, one pass.}
Because $\rhobar_i$ needs $\vrho_j$, hence $\vw_j$, hence $\vc_j$, hence $\nbr{j}$, a training step for cell $i$ reaches two graph hops even though the attention layer itself reaches one. Batching by \emph{contiguous spatial tile} makes this cheap, because the halo of extra cells around a tile is shared rather than disjoint. A tile's cells are the \emph{seeds}; their graph neighbours form \emph{ring one}, for which the context, the response posterior and the clean composition are all computed; the neighbours of ring one form \emph{ring two}, for which only $\vmu_z$ is computed to serve as attention sources. Both rings are built by look-ups in the same edge list that $\beta$ uses, never by a distance query: a cell across a pruned edge is not in ring one even if a radius search would return it, and if the halo and the leak kernel disagree the foreign influx is silently wrong. Attention edges are restricted to destinations in the seeds and ring one, since a ring-two cell's context would normalise over a truncated neighbourhood. The loss is evaluated on seeds only. With tiles of a few thousand cells the halo adds roughly a tenth to a fifth to the number of nodes processed per step, falling with tile size (measured, \cref{app:architecture}), cheaper than a cache of clean compositions that must be kept current; scattered seeds, by contrast, would cost many times more nodes than a contiguous tile, so spacing seeds apart is the expensive option, not the cheap one. The arithmetic of a pass is dominated by the encoder and decoder layers, at $O(NGh)$ for hidden width $h$; the foreign influx costs $O(|E|\,G)$ for $|E|$ directed edges, about $6NG$ here, though its gathered edge-by-gene rows are the largest intermediate tensor of a step. One consequence of tile batching should be stated: the ring of a training tile can contain cells of a held-out tile, whose counts then enter training as inputs under a stop-gradient. No gradient flows from them and no loss is evaluated on them, but they are used transductively \pending{quantify the fraction of training seeds with a held-out neighbour, or add a one-hop buffer around held-out tiles}.

\subsection{Model Selection and Diagnostics}
\label{sec:selection}

A model that explains everything with leakage reconstructs held-out cells well, so held-out reconstruction alone is not a valid stopping criterion. We stop on held-out reconstruction and on the agreement between $\vz$ and cell type \emph{jointly}: a checkpoint is accepted only if it improves held-out reconstruction while keeping the agreement (computed on a random subsample of all cells) at or above a fixed fraction ($0.9$) of its running maximum. Held-out cells are entire spatial tiles, never random cells, because spatial autocorrelation leaks the answer through neighbouring cells landing on both sides of a random split \citep{roberts2017}.

Five diagnostics accompany every fit. (i) The within-type coefficient of determination of $\vz_i$ regressed on $\vc_i$, against a control in which $\vc_i$ is permuted among cells of the same type; a pooled version would partly measure type separability, which both variables legitimately carry, and the within-type form is blind to it. This is the check on the value channel of the mirror described in \cref{sec:inference}. (ii) $\KL(q(\vw_i)\|p(\vw_i \mid \vc_i, t_i))$ per dimension. Its collapse to zero has three causes that the training curves alone cannot distinguish: posterior collapse under a large $\alphaw$, the mirror, and a response posterior starved of ego evidence. Each has its own check and fix (\cref{app:qw}), and the mirror must be ruled out before any free-bits style remedy \citep{kingma2016} is applied, since free bits would mask it and leave the anomaly score meaningless. (iii) The held-out probe of \cref{eq:probe}. (iv) The $\vz$--type agreement, the hard floor whose collapse is the marginal-invariance failure mode. (v) Two degeneracy checks in the opposite direction, asking whether $\vz$ has collapsed onto the label: the ratio $\I(\vz;t)/H(t)$ together with the within-type variance of $\vmu_z$, and the gap in held-out reconstruction between $\vz$ and a one-hot type code \pending{compute for the canonical fits}.

\subsection{What Is Identified, and the Leakage Sweep}
\label{sec:sweep}

\paragraph{Identifiability.}
Four statements can be made. The leak fraction $\kappa$ is not identified from the count matrix under the parametric forms we are willing to assume, and is swept. Given $\kappa$, the response coordinates are identified only up to a rotation: the isotropic prior of \cref{eq:gen-w} and the linear decoder are invariant under $\vw \mapsto \mathbf{Q}\vw$, $\vB \mapsto \vB\mathbf{Q}^\top$, $\mpsi \mapsto \mathbf{Q}\mpsi$ for orthogonal $\mathbf{Q}$ (the learned prior mean absorbs the last), while the fixed prior scale removes the rescaling freedom (\cref{sec:inference}); the diagonal posterior of \cref{eq:post-w} is not rotation-equivariant, so this is a symmetry of the model, not of the fitted objective. Through the softmax, each column of $\vB$ is further identified only up to an additive constant over genes, and so is $a(\vz)$ per cell; loadings are read only after centring each column over genes. The allocation between $\vz$ and $\vw$ is not pinned by the likelihood, since both posteriors see the counts; it rests on the conditional-invariance constraint, on the intrinsic path, and on $\dw < \dz$. For $\vw$ given the context this is in the spirit of the auxiliary-variable setting of \citet{khemakhem2020}, whose guarantee (recovery up to an affine map, under an injective additive-noise decoder and sufficient variability of the conditional prior) we do not establish for this likelihood; for $\vz$ no comparable guarantee exists without the constraint \citep{locatello2019}, and the constraint is what we verify with the probe rather than assume.

\paragraph{The sweep.}
The unit of inference is a sensitivity analysis with $\kappa$ as the sensitivity parameter, in the spirit of partial identification \citep{manski2003,rosenbaum2002}. The model is fitted independently at each $\kappa \in \{0, 0.05, 0.1, 0.2, 0.3, 0.4\}$ with several random seeds per grid point (three in the runs reported here, \cref{sec:sweep-results}), with every other weight held at the values calibrated at $\kappa = 0.1$ (\cref{app:calibration}) and the held-out probe checked at every $\kappa$ so that a change in a readout is not a change in the invariance. Every readout of interest, the loadings $\vB$, the per-type magnitude of the response, the spatial autocorrelation of each latent, the held-out probes and the agreement with external labels, is reported as a trajectory over $\kappa$ with its envelope over seeds. We call a readout stable when its sign and, for rankings, its order are unchanged over the grid and its seed envelope excludes zero at every grid point \todo{confirm this rule}. A stable readout is a finding; one that changes sign, vanishes, or is inconsistent across seeds is not separable from leakage on this data. Some readouts are expected to move with $\kappa$ by construction: as $\kappa$ grows, more of the resemblance between neighbours is assigned to the mixture and the fitted response magnitude shrinks, so a decline in $\|\vw\|$ along the grid is mechanical, whereas a change in the \emph{direction} of a programme or in which latent a target aligns with is not. Two consequences follow. Gene programmes are identified only up to the rotation and the per-column constants above, so raw loadings and probe directions are never compared across runs or seeds; only run-internal, rotation-invariant summaries are reported, and any gene-level statement is treated as relative to the run that produced it. And the reported range is not a confidence interval: it does not shrink with more cells, because the confound it expresses is not a sampling problem. An operating point, $\kappa = 0.1$ here, is used for figures and for the analyses that require a single fit, and is always reported alongside the sweep.

\paragraph{What $\vw$ is not.}
The context $\vc_i$ is the environment: deterministic, shared by co-located cells, a conditioning variable. If a niche descriptor were all that was wanted, one would drop the latent and feed $\vc_i$ to the decoder directly, which is roughly the node-centric model of \citet{fischer2023} and is simpler. The response $\vw_i$ is latent and per cell, and it earns its place only through within-niche heterogeneity: anomalies, per-cell responses, and counterfactuals with abduction. Whether it does so at a given operating point is an empirical question that the per-dimension divergence answers, and we report it as such (\cref{sec:w-degeneracy}).
